\documentclass[10pt,letterpaper]{article}
\usepackage{technical-implementation}
\setDocTitle{Timing Frame Rates and Exposure}
\setDocSubtitle{Technical reference | supplementary engineering with book crosswalk}
\setDocVersion{1.0}
\setDocStatus{Educational / Draft}
\setDocOwner{Game Development}
\setDocDate{2026-10-30}
\setDocSummary{Technical reference | supplementary engineering with book crosswalk}
\begin{document}
\makedoctitle
\makedocsummary
\tableofcontents\clearpage

\section{Scope and interpretation}
\textbf{Book foundation:} Chapters 9--17, 27--33, and 68--72 address production units, timing plans, exposures, holds, and synchronization. The following formulas and software requirements are supplementary engineering practice.
\section{Time bases}
Let $F=n/d$ frames/s (positive coprime integers), $\Delta t=d/n$ seconds/frame, and $t_k=kd/n$. Event boundaries must be represented in one explicitly named time base, ideally integer ticks or a rational pair. A frame index is not a duration.
\[T = \frac{N}{F},\qquad N=\operatorname{round}_{\mathrm{policy}}(TF),\qquad t_k=\frac{k}{F}.\]
Choose and disclose floor, ceil, or nearest-even rounding; never silently assume an exact integer frame count for a duration that is not frame aligned.
\section{Display frames, key poses, and drawing exposures}
A rendered frame can reuse a drawing; a drawing exposed on twos lasts two display frames in a traditional 24 fps shot. An 18-frame shot animated on twos contains at most nine distinct sequential drawings if every exposure has length two, subject to boundaries and reused cels. \emph{Twos} and 12 fps real-time playback are not identical: camera motion, compositing, interpolation and output sampling can still occur every display frame.
\section{Rational broadcast rates and timecode}
$24000/1001\approx 23.9760$ and $30000/1001\approx29.9700$ are rational rates, not integers. Timecode labels must declare non-drop-frame or drop-frame conventions; drop-frame timecode skips some labels to remain aligned to clock time and does not delete actual video frames. Do not convert timecode strings into elapsed time by assuming every minute has the same number of frame labels under drop-frame rules.
\section{Exposure sheet schema}
Every X-sheet row has \texttt{scene\_id}, \texttt{display\_frame}, \texttt{time\_seconds}, \texttt{drawing\_id}, \texttt{exposure\_length}, \texttt{dialogue\_cue}, \texttt{camera\_cue}, and \texttt{notes}. Exposures form a gap-free partition of the scene's display frames. The same drawing may occur on nonconsecutive rows for holds or cycles. A cue has an explicit start and end boundary.
\section{Worked example}
For 2.5 seconds at 24 fps, $N=60$ display frames. For a uniform two-frame exposure, create 30 assignments, each with a positive exposure length of two. A 12-frame hold starting at frame 18 occupies display frames 18--29; the next event begins at frame 30. Record whether indices are zero-based and whether intervals are half-open.
\section{Validation rules}
\begin{enumerate}
\item Sum of exposure lengths equals declared shot frame count; all lengths are positive integers.
\item Every displayed frame maps to exactly one drawing per animation layer.
\item Synchronization events use half-open intervals to prevent double-counted boundaries.
\item Conversion regression tests cover 24, 25, 30, 60, $24000/1001$, and $30000/1001$ fps.
\item Evaluators distinguish display cadence from simulation update cadence and authored-key density.
\end{enumerate}
\section{Common failure modes}
Frame/time rounding drift; interpreting ``24 frames'' as 24 seconds; unintentional one-frame overlap at cuts; exposed drawing cadence confused with renderer FPS; animation clocks incorrectly tied to a variable update rate.\par
\textbf{Source crosswalk:} Chapters 9--17, 27--33 and 68--72 are suggested readings; the validation schema and implementation guidance are independently developed.
\end{document}
